\chapter{Lecture 01}

\hfill {\footnotesize \textit{Date: 02/02/2026}}

Mathematical language of this course is linear algebra (of real vector spaces) and multivariable calculus. In this lecture, we will give a brief review of linear algebra.

\section{Review on linear algebra}

\textbf{Notation.} Let \( a_1, a_2, \ldots, a_n \) be a finite sequence of real numbers. We will denote the sum

\begin{equation}
	S = a_1 + a_2 + \cdots + a_n
	\label{eq:lecture-01-1}
\end{equation}

by

\begin{equation}
	S = \sum_{k=1}^{n} a_k.
	\label{eq:lecture-01-2}
\end{equation}

We have the following useful identities:

\begin{equation}
	\sum_{k=1}^{n} k = \frac{n(n+1)}{2},
	\label{eq:lecture-01-3}
\end{equation}

\begin{equation}
	\sum_{k=1}^{n} k^2 = \frac{n(n+1)(2n+1)}{6}.
	\label{eq:lecture-01-4}
\end{equation}

\subsection{Vectors}
\begin{definition}[\( \mathbb{R} \)-vector space]
	Let $V$ be a non-empty set. The set $V$ equipped with two binary operations
	\[
		+: V \times V \rightarrow V
	\]

	called \textit{vector addition}, and

	\[
		\cdot : \mathbb{R} \times V \rightarrow V
	\]

	called \textit{scalar multiplication} is called a \textbf{vector space over \( \mathbb{R} \)} or \( \mathbb{R} \)-\textbf{vector space} if for all $u,v,w \in V$ and all $a,b \in \mathbb{F}$, the following axioms hold:

	\begin{enumerate}
		\item \textbf{Closure under addition:}
		      \begin{equation}
			      u+v \in V
			      \label{eq:r-vector-space-1}
		      \end{equation}

		\item \textbf{Commutativity of addition:}
		      \begin{equation}
			      u+v=v+u
			      \label{eq:r-vector-space-2}
		      \end{equation}

		\item \textbf{Associativity of addition:}
		      \begin{equation}
			      (u+v)+w=u+(v+w)
			      \label{eq:r-vector-space-3}
		      \end{equation}

		\item \textbf{Existence of zero vector:}
		      There exists a vector $0 \in V$ such that
		      \begin{equation}
			      u+0=u
			      \label{eq:r-vector-space-4}
		      \end{equation}

		\item \textbf{Existence of additive inverse:}
		      For every $u\in V$, there exists $-u\in V$ such that
		      \begin{equation}
			      u+(-u)=0
			      \label{eq:r-vector-space-5}
		      \end{equation}

		\item \textbf{Closure under scalar multiplication:}
		      \begin{equation}
			      au\in V
			      \label{eq:r-vector-space-6}
		      \end{equation}

		\item \textbf{Distributivity of scalar multiplication over vector addition:}
		      \begin{equation}
			      a(u+v)=au+av
			      \label{eq:r-vector-space-7}
		      \end{equation}

		\item \textbf{Distributivity of scalar multiplication over field addition:}
		      \begin{equation}
			      (a+b)u=au+bu
			      \label{eq:r-vector-space-8}
		      \end{equation}

		\item \textbf{Associativity of scalar multiplication:}
		      \begin{equation}
			      a(bu)=(ab)u
			      \label{eq:r-vector-space-9}
		      \end{equation}

		\item \textbf{Identity element of scalar multiplication:}
		      \begin{equation}
			      1u=u
			      \label{eq:r-vector-space-10}
		      \end{equation}

	\end{enumerate}

	Hence, $(V,+,\cdot)$ is called a vector space over the field $\mathbb{F}$. Elements of \( V \) are called \textbf{vectors} and elements of \( \mathbb{R} \) are called \textbf{scalars}.
\end{definition}

\begin{example}

\end{example}

\begin{theorem}
	Blah
\end{theorem}

\subsection{Linear transformations}
